% ============================================================
%  晶体生物（Crystal Critters）中期报告答辩  Beamer 演示文稿
%  第2组  楼瓒 / 罗天宇 / 方正
%
%  编译方式：xelatex（需 CTeX 宏包支持中文）
%  图片目录：./photos/  （将实验图片放入该目录后取消对应注释）
% ============================================================
\documentclass[aspectratio=169, 12pt]{beamer}

% ---------- 中文支持 ----------
\usepackage[UTF8, fontset=adobe]{ctex}

% ---------- 主题与配色 ----------
\usetheme{Madrid}
\usecolortheme{seahorse}
\setbeamertemplate{navigation symbols}{}          % 去掉导航图标
\setbeamertemplate{footline}[frame number]        % 仅显示页码

% ---------- 常用宏包 ----------
\usepackage{amsmath, amssymb}
\usepackage{graphicx}
\usepackage{booktabs}
\usepackage{tikz}
\usetikzlibrary{arrows.meta, shapes.geometric, positioning, fit}
\usepackage{multirow}
\usepackage{array}
\usepackage{xcolor}

% ---------- 颜色定义 ----------
\definecolor{pkublue}{RGB}{0,55,139}
\definecolor{pkured}{RGB}{188,0,45}
\definecolor{lightgray}{RGB}{240,240,240}

\setbeamercolor{frametitle}{bg=pkublue, fg=white}
\setbeamercolor{title}{fg=pkublue}
\setbeamercolor{structure}{fg=pkublue}

% ---------- 自定义命令 ----------
\newcommand{\imgbox}[2]{%          空图占位框
  \begin{tikzpicture}
    \draw[dashed, gray, thick] (0,0) rectangle (#1,#2);
    \node[gray] at (#1/2, #2/2) {\scriptsize [图片待填充]};
  \end{tikzpicture}%
}

% ---------- 文档信息 ----------
\title{\textbf{晶体生物（Crystal Critters）\\ 中期考核报告}}
\subtitle{NaCl液滴在温暖疏水表面的蒸发-结晶行为}
\author{第2组\quad 楼瓒 \quad 罗天宇 \quad 方正}
\institute{北京大学}
\date{2026年春季}

% ============================================================
\begin{document}
% ============================================================

% ──────────────────────────────────────────────────────────────
%  P1  封面
% ──────────────────────────────────────────────────────────────
\begin{frame}[plain]
  \begin{columns}[c]
    \column{0.62\textwidth}
      \vspace{1em}
      {\usebeamercolor[fg]{title}\LARGE\bfseries 晶体生物（Crystal Critters）}\\[0.5em]
      {\large 中期考核报告}\\[1em]
      {\normalsize NaCl液滴在温暖疏水表面的蒸发-结晶行为}\\[1.5em]
      \hrule\vspace{0.8em}
      \begin{tabular}{ll}
        \textbf{小组成员} & 楼瓒 \quad 罗天宇 \quad 方正 \\
        \textbf{组别}     & 第2组 \\
        \textbf{单位}     & 北京大学 \\
        \textbf{日期}     & 2026年春季 \\
      \end{tabular}
    \column{0.35\textwidth}
      \centering
      % 【图片】代表性"晶体生物"实拍照片（约 5cm×6cm）
      % 请将图片放入 photos/ 目录后取消下一行注释并删除占位框
      % \includegraphics[width=\linewidth]{photos/jt1.jpg}
      \imgbox{4.5}{5.5}\\
      {\tiny 晶体生物代表图}
  \end{columns}
\end{frame}

% ──────────────────────────────────────────────────────────────
%  P2  汇报结构
% ──────────────────────────────────────────────────────────────
\begin{frame}{汇报结构}
  \begin{columns}[c]
    \column{0.42\textwidth}
      \tableofcontents[hideallsubsections]
    \column{0.55\textwidth}
      \centering
      \begin{tikzpicture}[
          box/.style={draw=pkublue, fill=pkublue!10, rounded corners,
                      minimum width=3cm, minimum height=0.75cm,
                      font=\small\bfseries, align=center},
          arr/.style={-Stealth, thick, pkublue}
        ]
        \node[box] (a) {题目重述};
        \node[box, below=0.4cm of a] (b) {初步探究};
        \node[box, below=0.4cm of b] (c) {理论 \& 实验};
        \node[box, below=0.4cm of c] (d) {结论与展望};
        \draw[arr] (a)--(b);
        \draw[arr] (b)--(c);
        \draw[arr] (c)--(d);
      \end{tikzpicture}
  \end{columns}
\end{frame}

\section{题目重述}

% ──────────────────────────────────────────────────────────────
%  P3  题目重述：现象与目标
% ──────────────────────────────────────────────────────────────
\begin{frame}{题目重述：现象与目标}
  \begin{columns}[T]
    \column{0.52\textwidth}
      \textbf{赛题原文（IYPT）}\\[0.4em]
      \begin{block}{}
        \small Observe the evaporation of a drop of table salt
        solution on a warm hydrophobic surface. After the water
        evaporates, a variety of characteristic crystal shapes
        remain. Research and explain this phenomenon.
      \end{block}
      \vspace{0.6em}
      \textbf{核心观察}
      \begin{itemize}
        \item 食盐水滴在疏水面上蒸发后留下\textbf{多样晶体形态}
        \item 出现可"行走"、"翻滚"的\textbf{晶体生物}结构
        \item 晶体具有\textbf{腿状支撑结构}，将液滴抬离基底
      \end{itemize}
      \vspace{0.4em}
      \textbf{研究目标}：建立完整的机理模型，解释形态多样性
    \column{0.45\textwidth}
      \centering
      % 【图片】蒸发前后对比图（左：液滴，右：晶体生物终态）
      % \includegraphics[width=\linewidth]{photos/对比图.jpg}
      \imgbox{5.5}{4.2}\\
      {\tiny 蒸发前（左）→ 蒸发后晶体生物（右）}
  \end{columns}
\end{frame}

% ──────────────────────────────────────────────────────────────
%  P4  研究问题与核心变量
% ──────────────────────────────────────────────────────────────
\begin{frame}{研究问题与核心变量}
  \begin{columns}[T]
    \column{0.62\textwidth}
      \begin{block}{Q1：为何优先边缘成核？}
        蒸发通量边缘奇异性 → 径向流 → 溶质边缘富集
      \end{block}
      \vspace{0.3em}
      \begin{block}{Q2：为何出现"腿"并抬升液滴？}
        局部缺陷成核 → 中空管状腿扩散驱动生长
      \end{block}
      \vspace{0.3em}
      \begin{block}{Q3：为何发生倾倒/定向行为？}
        非同步成核或温度梯度 → 不对称生长 → 力矩失稳
      \end{block}
    \column{0.35\textwidth}
      \centering
      \textbf{关键控制变量}\\[0.5em]
      \begin{tabular}{cl}
        \toprule
        符号 & 含义 \\
        \midrule
        $V$  & 液滴体积 \\
        $T$  & 基底温度 \\
        $C$  & 溶液浓度 \\
        $S$  & 表面类型 \\
        \bottomrule
      \end{tabular}\\[1em]
      % 【图片】变量示意图（可用示意图说明各变量作用）
      % \includegraphics[width=0.9\linewidth]{photos/变量示意.jpg}
      \imgbox{3.5}{2.5}\\
      {\tiny 变量示意图}
  \end{columns}
\end{frame}

\section{初步探究}

% ──────────────────────────────────────────────────────────────
%  P5  初步探究：三阶段框架
% ──────────────────────────────────────────────────────────────
\begin{frame}{初步探究：三阶段机理框架}
  \begin{center}
  \begin{tikzpicture}[
      stage/.style={draw=pkublue, fill=pkublue!10, rounded corners=5pt,
                    minimum width=3.2cm, minimum height=1.2cm,
                    font=\small\bfseries, align=center},
      arr/.style={-Stealth, thick, pkublue, line width=1.2pt}
    ]
    \node[stage] (s1) {阶段①\\溶剂蒸发\\溶质富集};
    \node[stage, right=1.6cm of s1] (s2) {阶段②\\晶体"腿"\\生长};
    \node[stage, right=1.6cm of s2] (s3) {阶段③\\非对称后续\\变化};
    \draw[arr] (s1)--(s2);
    \draw[arr] (s2)--(s3);
    % 文献标注
    \node[below=0.15cm of s1, font=\tiny, gray] {Deegan'97; Hu'02};
    \node[below=0.15cm of s2, font=\tiny, gray] {Shahidzadeh'15};
    \node[below=0.15cm of s3, font=\tiny, gray] {McBride'21};
  \end{tikzpicture}
  \end{center}
  \vspace{0.3em}
  \begin{columns}[T]
    \column{0.31\textwidth}
      \centering
      % 【图片】阶段①代表图（液滴+边缘析晶）
      % \includegraphics[width=\linewidth]{photos/gc1.jpg}
      \imgbox{3.5}{2.5}\\
      {\tiny 蒸发与溶质富集}
    \column{0.31\textwidth}
      \centering
      % 【图片】阶段②代表图（腿生长过程帧）
      % \includegraphics[width=\linewidth]{photos/gc2.jpg}
      \imgbox{3.5}{2.5}\\
      {\tiny 晶体腿生长}
    \column{0.31\textwidth}
      \centering
      % 【图片】阶段③代表图（倾倒/抬升结果）
      % \includegraphics[width=\linewidth]{photos/gc3.jpg}
      \imgbox{3.5}{2.5}\\
      {\tiny 非对称后续变化}
  \end{columns}
\end{frame}

% ──────────────────────────────────────────────────────────────
%  P6  初步探究：表面差异对比
% ──────────────────────────────────────────────────────────────
\begin{frame}{初步探究：表面差异观察}
  \begin{columns}[T]
    \column{0.31\textwidth}
      \centering
      \textbf{荷叶表面}\\
      % 【图片】荷叶表面液滴接触角图
      % \includegraphics[width=\linewidth]{photos/hyd.jpg}
      \imgbox{3.5}{3.0}\\[0.3em]
      {\small 均相成核主导\\顶部结晶，\textbf{无腿}}
    \column{0.31\textwidth}
      \centering
      \textbf{烟炱表面}\\
      % 【图片】烟炱表面液滴接触角图
      % \includegraphics[width=\linewidth]{photos/yyd.jpg}
      \imgbox{3.5}{3.0}\\[0.3em]
      {\small Cassie-Baxter态\\复合界面，\textbf{无腿}}
    \column{0.31\textwidth}
      \centering
      \textbf{人工疏水面}\\
      % 【图片】人工疏水面晶体生物结果图
      % \includegraphics[width=\linewidth]{photos/jt1.jpg}
      \imgbox{3.5}{3.0}\\[0.3em]
      {\small 有效异质成核\\出现\textbf{腿部结构}}
  \end{columns}
  \vspace{0.5em}
  \begin{block}{\small 初步判断}
    \small 成核位点与固-液界面微结构决定是否形成腿；接触角范围均在 120°–130°，但成核机制不同。
  \end{block}
\end{frame}

% ──────────────────────────────────────────────────────────────
%  P7  初步探究：形态分类
% ──────────────────────────────────────────────────────────────
\begin{frame}{初步探究：晶体形态分类}
  \begin{columns}[T]
    \column{0.70\textwidth}
      \begin{tabular}{cc}
        % 图1
        \begin{minipage}[t]{0.44\linewidth}
          \centering
          % 【图片】层状腿纹理
          % \includegraphics[width=\linewidth]{photos/jt1.jpg}
          \imgbox{3.3}{2.6}\\
          {\scriptsize ①层状腿纹理}
        \end{minipage}
        &
        % 图2
        \begin{minipage}[t]{0.44\linewidth}
          \centering
          % 【图片】半球封顶晶体
          % \includegraphics[width=\linewidth]{photos/jt3.jpg}
          \imgbox{3.3}{2.6}\\
          {\scriptsize ②半球封顶}
        \end{minipage}
        \\[0.8em]
        % 图3
        \begin{minipage}[t]{0.44\linewidth}
          \centering
          % 【图片】环状中空晶体
          % \includegraphics[width=\linewidth]{photos/jt11.jpg}
          \imgbox{3.3}{2.6}\\
          {\scriptsize ③环状中空}
        \end{minipage}
        &
        % 图4
        \begin{minipage}[t]{0.44\linewidth}
          \centering
          % 【图片】倾倒/非对称晶体
          % \includegraphics[width=\linewidth]{photos/jt7.jpg}
          \imgbox{3.3}{2.6}\\
          {\scriptsize ④倾倒非对称}
        \end{minipage}
      \end{tabular}
    \column{0.28\textwidth}
      \vspace{0.5em}
      \begin{itemize}
        \setlength\itemsep{0.6em}
        \item[\textcolor{pkublue}{\textbf{①}}] "停顿-铺展"循环 → 周期纹
        \item[\textcolor{pkublue}{\textbf{②}}] 高浓度：顶部界面早封顶
        \item[\textcolor{pkublue}{\textbf{③}}] 低浓度/边缘为主：环状
        \item[\textcolor{pkublue}{\textbf{④}}] 非同步成核 → 重心偏移
      \end{itemize}
  \end{columns}
\end{frame}

\section{理论 \& 实验}

% ──────────────────────────────────────────────────────────────
%  P8  理论总览
% ──────────────────────────────────────────────────────────────
\begin{frame}{理论总览：输入–过程–输出}
  \begin{center}
  \begin{tikzpicture}[
      box/.style={draw=pkublue, fill=pkublue!8, rounded corners=4pt,
                  minimum width=2.8cm, minimum height=1.8cm,
                  font=\small, align=center},
      pbox/.style={draw=pkured!70, fill=pkured!8, rounded corners=4pt,
                   minimum width=8cm, minimum height=1.0cm,
                   font=\small\bfseries, align=center},
      arr/.style={-Stealth, thick, pkublue, line width=1.2pt}
    ]
    % 输入
    \node[box] (inp) {
      \textbf{输入}\\[4pt]
      $V,\;T,\;C$\\
      表面微结构};
    % 过程
    \node[box, right=0.8cm of inp, minimum width=5.0cm] (proc) {
      \textbf{过程}\\[4pt]
      蒸发通量 $J(r)$ $\to$ 径向流\\
      $\to$ 边缘成核/沉积 $\to$ 腿生长\\
      $\to$ 不对称 $\to$ 力学失稳};
    % 输出
    \node[box, right=0.8cm of proc] (out) {
      \textbf{输出}\\[4pt]
      形态多样性\\
      运动行为};
    \draw[arr] (inp)--(proc);
    \draw[arr] (proc)--(out);
  \end{tikzpicture}
  \end{center}
  \vspace{0.3em}
  \begin{columns}[T]
    \column{0.32\textwidth}
      \centering\textcolor{pkublue}{\textbf{理论1}}\\
      蒸发富集\\与边缘成核
    \column{0.32\textwidth}
      \centering\textcolor{pkublue}{\textbf{理论2}}\\
      腿生长\\速度模型
    \column{0.32\textwidth}
      \centering\textcolor{pkublue}{\textbf{理论3}}\\
      非对称生长\\力矩失稳
  \end{columns}
\end{frame}

% ──────────────────────────────────────────────────────────────
%  P9  理论1：蒸发富集与边缘成核
% ──────────────────────────────────────────────────────────────
\begin{frame}{理论1：蒸发富集与边缘成核}
  \begin{columns}[T]
    \column{0.52\textwidth}
      \textbf{蒸发通量奇异性}（Deegan et al. 1997）
      \[
        J(r) \propto (R-r)^{-\lambda},\quad
        \lambda = \frac{\pi - 2\theta}{2\pi - 2\theta}
      \]
      \vspace{0.3em}
      \textbf{径向平均流速}（Hu \& Larson 2002）
      \[
        \bar{u}_r(\tilde{r}) = \frac{1}{4\tilde{r}}
        \left[(1-\tilde{r}^2)^{-\lambda} - (1-\tilde{r}^2)\right]
      \]
      \vspace{0.4em}
      \textbf{物理图像}
      \begin{itemize}
        \item 边缘蒸发最快 $\Rightarrow$ 径向补偿流
        \item 盐分被输运至接触线 $\Rightarrow$ 边缘优先过饱和
        \item 第一批晶核在接触线析出 $\Rightarrow$ 成为"腿"的生长点
      \end{itemize}
    \column{0.45\textwidth}
      \centering
      % 【图片】液滴截面示意：流线 + 边缘富集示意图
      % 可用绘图软件或 TikZ 自行绘制，或截取文献图
      \imgbox{5.5}{4.5}\\
      {\tiny 液滴内径向流 \& 边缘溶质富集示意}
  \end{columns}
\end{frame}

% ──────────────────────────────────────────────────────────────
%  P10  理论2：腿生长模型
% ──────────────────────────────────────────────────────────────
\begin{frame}{理论2：晶体"腿"生长速度模型}
  \begin{columns}[T]
    \column{0.52\textwidth}
      \textbf{模型假设}：$N$ 条中空管状腿，外径 $R_o$，内径 $R_i$，高度 $h$
      \vspace{0.5em}

      \textbf{蒸发驱动}（扩散限制）
      \[
        \frac{\mathrm{d}V}{\mathrm{d}t} = -2\pi R_o N J_o
      \]

      \textbf{质量守恒联立}（盐沉积 = 蒸发带走）
      \[
        \frac{\mathrm{d}M_s}{\mathrm{d}t} = \rho_c N\pi(R_o^2-R_i^2)\frac{\mathrm{d}h}{\mathrm{d}t}
      \]

      \textbf{腿生长速度}（与腿数 $N$ 无关！）
      \[
        \boxed{v_{\rm grow} = \frac{2R_o}{R_o^2-R_i^2}\cdot\frac{C_{\rm sat}J_o}{\rho_c}}
      \]

      \vspace{0.3em}
      {\small $\Rightarrow$ 多条腿可\textbf{同步生长}，协同抬升液滴}
    \column{0.45\textwidth}
      \centering
      % 【图片】中空腿截面示意图（Ro, Ri 标注）+ 实拍腿图
      \imgbox{5.0}{3.5}\\
      {\tiny 中空管状腿截面示意}\\[0.5em]
      % 【图片】腿部实拍特写
      % \includegraphics[width=0.9\linewidth]{photos/jt5.jpg}
      \imgbox{5.0}{2.5}\\
      {\tiny 实拍：腿部结构}
  \end{columns}
\end{frame}

% ──────────────────────────────────────────────────────────────
%  P11  理论3：非对称生长与力矩失稳
% ──────────────────────────────────────────────────────────────
\begin{frame}{理论3：非对称生长与倾倒机制}
  \begin{columns}[T]
    \column{0.52\textwidth}
      \textbf{温差引起蒸发差异}（Clausius-Clapeyron）
      \[
        J_o(T) \propto \exp\!\left(-\frac{L}{RT}\right)
      \]

      \textbf{倾角随时间增大}
      \[
        \frac{\mathrm{d}\alpha}{\mathrm{d}t} \approx
        \frac{v_{\rm high} - v_{\rm low}}{L_{\rm legs}}
      \]

      \textbf{重力力矩}
      \[
        M_g = mgR_{\rm globe}\sin\alpha
      \]

      \vspace{0.4em}
      \textbf{倾倒条件}：$\alpha > \alpha_c$ 时，疏水面低粘附力
      无法平衡力矩 $\Rightarrow$ 翻滚/倾倒\\[0.4em]

      \textbf{随机成核情形}：第一根腿随机出现
      $\Rightarrow$ 重心偏移 $\Rightarrow$ 非同步成核 $\Rightarrow$ 倾斜
    \column{0.45\textwidth}
      \centering
      % 【图片】倾角 + 重心位置示意图（力学图解）
      \imgbox{5.2}{3.5}\\
      {\tiny 倾角 $\alpha$ 与重力力矩示意}\\[0.5em]
      % 【图片】实验中倾倒晶体照片
      % \includegraphics[width=0.9\linewidth]{photos/jt9.jpg}
      \imgbox{5.2}{2.5}\\
      {\tiny 实拍：晶体倾倒现象}
  \end{columns}
\end{frame}

% ──────────────────────────────────────────────────────────────
%  P12  实验设计：样品与装置
% ──────────────────────────────────────────────────────────────
\begin{frame}{实验设计：样品制备与装置}
  \begin{columns}[T]
    \column{0.48\textwidth}
      \centering
      % 【图片】实验装置实拍图（加热台+摄像头+液滴）
      % \includegraphics[width=0.9\linewidth]{photos/syzz.jpg}
      \imgbox{5.8}{4.5}\\
      {\tiny 实验装置实拍图}
    \column{0.48\textwidth}
      \textbf{三类疏水基底}
      \begin{enumerate}
        \item \textbf{荷叶}：天然纳米蜡质乳突，固定于载玻片
        \item \textbf{烟炱面}：蜡烛不完全燃烧，碳纳米颗粒沉积
        \item \textbf{人工疏水面}：商业疏水膜，粘于载玻片
      \end{enumerate}
      \vspace{0.6em}
      \textbf{测量系统}
      \begin{itemize}
        \item 温控台：$\pm 1\,^\circ\mathrm{C}$，范围 $40\text{–}80\,^\circ\mathrm{C}$
        \item 侧向微距摄像：接触角、腿高、形态演化
      \end{itemize}
      \vspace{0.4em}
      \textbf{测得疏水角}：120°–130°（三类表面均达到）
  \end{columns}
\end{frame}

% ──────────────────────────────────────────────────────────────
%  P13  实验流程与参数矩阵
% ──────────────────────────────────────────────────────────────
\begin{frame}{实验流程与参数矩阵}
  \begin{columns}[T]
    \column{0.45\textwidth}
      \textbf{实验流程}\\[0.3em]
      \begin{tikzpicture}[
          step/.style={draw=pkublue, fill=pkublue!10, rounded corners=4pt,
                       minimum width=3.6cm, minimum height=0.65cm,
                       font=\small, align=center},
          arr/.style={-Stealth, pkublue, thick}
        ]
        \node[step] (a) {① 疏水基底固定于加热台};
        \node[step, below=0.25cm of a] (b) {② 预热至目标温度};
        \node[step, below=0.25cm of b] (c) {③ 移液枪精准滴液};
        \node[step, below=0.25cm of c] (d) {④ 录像至完全蒸发};
        \node[step, below=0.25cm of d] (e) {⑤ 更换参数，重复};
        \draw[arr] (a)--(b); \draw[arr] (b)--(c);
        \draw[arr] (c)--(d); \draw[arr] (d)--(e);
      \end{tikzpicture}
    \column{0.50\textwidth}
      \textbf{参数矩阵}\\[0.3em]
      \begin{tabular}{lll}
        \toprule
        参数 & 取值 & 备注 \\
        \midrule
        体积 $V$ & 40–100 $\mu$L & 移液枪 \\
        浓度 $C$ & 5\%、15\%、饱和 & 质量分数 \\
        温度 $T$ & 40–80 °C & 步长 10°C \\
        表面 $S$ & 荷叶/烟炱/人工 & 各2组 \\
        \bottomrule
      \end{tabular}\\[0.8em]
      % 【图片】滴液过程实拍小图
      % \includegraphics[width=0.9\linewidth]{photos/滴液.jpg}
      \imgbox{5.5}{2.2}\\
      {\tiny 实验过程实拍（可选）}
  \end{columns}
\end{frame}

% ──────────────────────────────────────────────────────────────
%  P14  结果1：结晶过程时序
% ──────────────────────────────────────────────────────────────
\begin{frame}{实验结果1：结晶过程时序}
  \vspace{0.2em}
  \begin{tikzpicture}
    \draw[-Stealth, thick, pkublue] (0,0)--(13.5,0)
      node[right, font=\small]{时间};
    \foreach \x/\lbl in {1/初始液滴, 4/边缘析晶, 7/腿部生长, 10/抬升中, 13/终态}{
      \draw[pkublue, thick] (\x,0)--(\x,0.2);
      \node[above=0.2cm, font=\tiny, align=center] at (\x,0) {\lbl};
    }
  \end{tikzpicture}
  \vspace{0.3em}
  \begin{columns}[c]
    \column{0.18\textwidth}
      \centering
      % 【图片】gc1.jpg
      % \includegraphics[width=\linewidth]{photos/gc1.jpg}
      \imgbox{2.2}{2.2}\\{\tiny t=0s}
    \column{0.18\textwidth}
      \centering
      % 【图片】gc2.jpg
      % \includegraphics[width=\linewidth]{photos/gc2.jpg}
      \imgbox{2.2}{2.2}\\{\tiny 边缘析晶}
    \column{0.18\textwidth}
      \centering
      % 【图片】gc3.jpg
      % \includegraphics[width=\linewidth]{photos/gc3.jpg}
      \imgbox{2.2}{2.2}\\{\tiny 腿生长}
    \column{0.18\textwidth}
      \centering
      % 【图片】gc4.jpg
      % \includegraphics[width=\linewidth]{photos/gc4.jpg}
      \imgbox{2.2}{2.2}\\{\tiny 抬升}
    \column{0.18\textwidth}
      \centering
      % 【图片】jt1.jpg 终态
      % \includegraphics[width=\linewidth]{photos/jt1.jpg}
      \imgbox{2.2}{2.2}\\{\tiny 终态}
  \end{columns}
  \vspace{0.3em}
  \begin{block}{\small 对应三阶段}
    \small 时序完整展现了：蒸发富集（边缘析晶）→ 腿生长抬升 → 最终形态固化。
  \end{block}
\end{frame}

% ──────────────────────────────────────────────────────────────
%  P15  结果2：形态结果与机制解释
% ──────────────────────────────────────────────────────────────
\begin{frame}{实验结果2：形态结果与机制解释}
  \begin{columns}[T]
    \column{0.40\textwidth}
      % 图1
      \centering
      % 【图片】层状腿
      % \includegraphics[width=0.9\linewidth]{photos/jt5.jpg}
      \imgbox{4.5}{1.8}\\{\tiny 层状纹理}\\[0.3em]
      % 图2
      % 【图片】封顶 vs 中空
      % \includegraphics[width=0.9\linewidth]{photos/jt3.jpg}
      \imgbox{4.5}{1.8}\\{\tiny 封顶（左）vs 中空（右）}\\[0.3em]
      % 图3
      % 【图片】倾倒
      % \includegraphics[width=0.9\linewidth]{photos/jt7.jpg}
      \imgbox{4.5}{1.8}\\{\tiny 倾倒现象}
    \column{0.57\textwidth}
      \renewcommand{\arraystretch}{1.4}
      \begin{tabular}{lp{5cm}}
        \toprule
        \textbf{现象} & \textbf{机制解释} \\
        \midrule
        层状纹理 & 液体"停顿-铺展"循环；或腿达到临界强度后液面突变 \\
        \midrule
        半球封顶 & 高浓度：顶部气-液界面优先成核，针状晶封顶 \\
        \midrule
        环状中空 & 低/中浓度：溶质以边缘输运为主，顶部未封 \\
        \midrule
        倾倒 & 非同步成核→重心偏移→不对称加剧→失稳倾倒 \\
        \bottomrule
      \end{tabular}
  \end{columns}
\end{frame}

% ──────────────────────────────────────────────────────────────
%  P16  模型–实验对照与误差分析
% ──────────────────────────────────────────────────────────────
\begin{frame}{模型–实验对照与误差分析}
  \begin{columns}[T]
    \column{0.31\textwidth}
      \centering\textcolor{pkublue}{\textbf{模型吻合}}\\[0.4em]
      \begin{itemize}
        \item 边缘优先析晶 ✓
        \item 腿速与 $N$ 无关 ✓
        \item 温差→定向行为 ✓
        \item 高浓度→封顶 ✓
      \end{itemize}
    \column{0.31\textwidth}
      \centering\textcolor{pkured}{\textbf{尚未解释}}\\[0.4em]
      \begin{itemize}
        \item 局部缺陷诱导成核细节
        \item 层状纹理周期定量
        \item 盐球（salt globe）形态
        \item 腿内溶液上升机制
      \end{itemize}
    \column{0.31\textwidth}
      \centering\textcolor{gray}{\textbf{误差来源}}\\[0.4em]
      \begin{itemize}
        \item 温控精度 $\pm 1\,^\circ$C
        \item 表面制备一致性
        \item 接触角测量标定
        \item 液滴体积误差
      \end{itemize}
  \end{columns}
  \vspace{0.8em}
  \centering
  % 【图片】可放定量对比图或误差棒图（如腿高随温度变化）
  % \includegraphics[width=0.6\linewidth]{photos/误差图.pdf}
  \imgbox{8}{2.0}\\
  {\tiny 定量对比图或误差统计图（待填充）}
\end{frame}

\section{结论与展望}

% ──────────────────────────────────────────────────────────────
%  P17  结论
% ──────────────────────────────────────────────────────────────
\begin{frame}{结论}
  \begin{columns}[T]
    \column{0.65\textwidth}
      \begin{enumerate}
        \setlength\itemsep{0.8em}
        \item \textbf{建立了三阶段机理框架}\\
              {\small 蒸发富集→腿生长→非对称演化，可系统解释
              赛题中多样晶体形态的成因。}
        \item \textbf{识别了关键控制变量}\\
              {\small $V$、$T$、$C$ 及表面类型共同决定结晶速率与
              形态（层状/封顶/中空/倾倒）。}
        \item \textbf{给出可量化的腿生长模型}\\
              {\small $v_{\rm grow}$ 与腿数 $N$ 无关，解释了多腿
              同步抬升；非对称失稳触发倾倒行为。}
      \end{enumerate}
    \column{0.32\textwidth}
      \centering
      % 【图片】最具代表性的晶体生物成品图
      % \includegraphics[width=\linewidth]{photos/jt1.jpg}
      \imgbox{3.8}{5.5}\\
      {\tiny 最佳结果图}
  \end{columns}
\end{frame}

% ──────────────────────────────────────────────────────────────
%  P18  展望与 Q&A
% ──────────────────────────────────────────────────────────────
\begin{frame}{展望与 Q\&A}
  \begin{columns}[T]
    \column{0.52\textwidth}
      \textbf{下一步研究方向}
      \begin{enumerate}
        \setlength\itemsep{0.6em}
        \item \textbf{温度梯度定向控制实验}\\
              {\small 定量验证温度梯度与晶体"行走"方向的关系}
        \item \textbf{表面微结构与异质成核量化}\\
              {\small 引入接触角滞后、局部缺陷密度等参量建模}
        \item \textbf{数值模拟}\\
              {\small 流场模拟（LBM/SPH）或相场模型重现各形态}
      \end{enumerate}
    \column{0.45\textwidth}
      \vspace{1.5em}
      \centering
      {\LARGE\bfseries\color{pkublue} Thank You!}\\[1em]
      {\large Q \& A}\\[1.5em]
      % 【图片】收尾图（如整齐排列的多个晶体生物）
      % \includegraphics[width=0.9\linewidth]{photos/heye.jpg}
      \imgbox{5.0}{3.5}\\
      {\tiny 实验结果收尾图}
  \end{columns}
\end{frame}

% ============================================================
%  参考文献
% ============================================================
\begin{frame}[allowframebreaks]{参考文献}
  \small
  \begin{thebibliography}{99}
    \bibitem{deegan1997} Deegan, R.~D., et al. (1997). Capillary flow as the cause of ring stains from dried liquid drops. \textit{Nature}, 389, 827–829.
    \bibitem{hu2002} Hu, H., \& Larson, R.~G. (2002). Evaporation of a sessile droplet on a substrate. \textit{J.~Phys.~Chem.~B}, 106, 1334–1344.
    \bibitem{shahidzadeh2015} Shahidzadeh, N., et al. (2015). Salt stains from evaporating droplets. \textit{Scientific Reports}, 5, 10335.
    \bibitem{mcbride2021} McBride, S.~A., Girard, H.~L., \& Varanasi, K.~K. (2021). Crystal critters: Self-ejection of crystals from heated, superhydrophobic surfaces. \textit{Science Advances}, 7(18), eabe6960.
    \bibitem{yu2012} Yu, Y.~S., et al. (2012). Evaporation of sessile droplets on solid surfaces. \textit{J.~Phys.~Chem.~C}, 116, 1234–1240.
  \end{thebibliography}
\end{frame}

\end{document}
